\section{Mathematical explanation and related systems}
\subsection{Error geometry and the fitting objective}
\label{app:geometry}
For a fixed library, Equation~\eqref{eq:riskmain} expresses expected squared
relative error as $R(w)=w^\top Cw$, with $C=\E[G(x,y)]$. The error product
matrix is uncentered, so shared bias remains part of the risk. Its diagonal
describes individual accuracy and its cross terms describe error alignment
\citep{breiman1996,krogh1994}. Writing $d_m=\widehat C_{mm}>0$, minimization
of the diagonal approximation $\sum_m d_mw_m^2$ under simplex constraints
gives $2d_mw_m=\lambda$. The scalar $\lambda$ enforces the unit sum constraint,
so normalization yields $w_m\propto d_m^{-1}$. This recovers the inverse error
mixture before its numerical floor, without assuming independent model errors.
If some $d_m=0$, an optimum is supported on the zero error entries.

For one case, let $e_m$ be model $m$'s normalized error field and
$r_m=\norm{e_m}_Q$ its magnitude. The improvement relative to errors of the
same sizes that all point in one direction is
\begin{equation}
 \left(\sum_m w_m r_m\right)^2-w^\top Gw
 =2\sum_{m<n}w_mw_n\left(r_mr_n-\langle e_m,e_n\rangle_Q\right)\geq0.
 \label{eq:direction}
\end{equation}
Here $G$ is the same case's error product matrix, and $m<n$ counts each model
pair once. Expansion gives the equality and Cauchy Schwarz gives nonnegativity.
The identity explains how errors in different locations or directions can help
either mixture. Negative correlation is not required, and the identity does
not establish that a small fitting set will discover useful weights.

The fitting criterion can change the preferred mixture \citep{acar2009,acar2015}.
Our joint fit minimizes mean squared relative error, whereas reporting uses
\begin{equation}
 L(w)=\E\!\left[\sqrt{w^\top G(x,y)w}\right].
 \label{eq:reportedrisk}
\end{equation}
Here $L$ averages the relative $L^2$ error over inputs and targets, taking
the square root separately for each field. Jensen's inequality gives
$L(w)\leq\sqrt{R(w)}$ for fixed weights, but does not equate their minimizers.
For example, two equally likely scalar cases with normalized model errors
$e_A=(0,2)$ and $e_B=(1.1,1.1)$ give
\begin{equation}
 L(t)=1.1-0.1t,\qquad R(t)=1.21-0.22t+1.01t^2,
\end{equation}
where $t\in[0,1]$ weights model A and $1-t$ weights B. The minimum of $L$
is at $t=1$, while $R$ is minimized at $t=11/101$, where $L=110/101$.
Thus fitting squared error can be about $8.9\%$ worse on the reported metric
even with unlimited fitting cases. Constant spatial fields give the same example.

\paragraph{Direct fitting of the reporting metric.}
A standard convex alternative fits the unsquared error directly:
\begin{equation}
 \min_{w\in\Delta_M}\frac1K\sum_{i=1}^K\sqrt{w^\top G_iw}
 =\min_{w\in\Delta_M}\frac1K\sum_{i=1}^K\norm{A_iw}_2,
 \qquad A_i^\top A_i=G_i.
 \label{eq:metricfit}
\end{equation}
Here $A_i$ factors the error matrix of fitting field $i$, $K$ counts fields,
and $\Delta_M$ contains nonnegative weights summing to one. Each term is a
norm of a linear map. Introducing upper bounds $t_i$ with
$\norm{A_iw}_2\leq t_i$ gives a second order cone problem. We evaluate this
loss control on 51 partitions of the original 17 PDE tasks at $K=5$, keeping
all predictions and case assignments fixed. The weights were saved before
their new scores were computed, but the predictions had previously been
inspected, making this a retrospective diagnostic separate from the three
main rules.

CVXPY 1.9.2 and Clarabel 0.11.1 solve the scaled problem with absolute and
relative gap tolerances $10^{-10}$ and feasibility tolerance $10^{-9}$.
All 51 solves returned optimal status. Eigenvalues below zero are clipped
only within $10^{-12}$ of spectral scale, with larger violations rejected.
The worst relative eigenvalue was approximately $-1.04\times10^{-18}$ and the
largest simplex residual before clipping and renormalization was
$4.15\times10^{-11}$. Source identities, reference error replay, and solver
records are retained with the locked fits.

\begin{table}[h!]
\centering\small
\caption{\textbf{Fitting the reporting metric directly.} Ratios compare with
the Selected model on the 17 PDE diagnostic tasks at five fitting examples.
Wins and losses require more than 1\% change after averaging three partitions.
The direct loss fit is a supplementary control, not a fourth main rule.}
\label{tab:losscontrol}\begin{tabular}{lrrrr}
\toprule
Method & Class ratio & Dataset ratio & Wins & Losses \\
\midrule
Inverse error mixture & 0.882 & 0.924 & 11 & 4 \\
Fitted mixture & 0.876 & 0.918 & 12 & 5 \\
Relative $L^2$ fit & 0.859 & 0.898 & 12 & 5 \\
\bottomrule
\end{tabular}

\end{table}

Direct relative $L^2$ fitting reduces class aggregated error by $2.0\%$ and
dataset aggregated error by $2.1\%$ relative to the original squared error
fit. Nine datasets improve by more than 1\%, none worsens by that margin,
and eight remain within 1\%. Against the Selected model, 12 improve and five
worsen. The result identifies a consequential loss choice without establishing
a new weighting principle or independent validation.

\subsection{What estimating the error matrix guarantees}
\label{app:proof}
\begin{proposition}[Error estimation and mixture risk]
\label{prop:bound}
If $\max_{m,n}|\widehat C_{mn}-C_{mn}|\leq\epsilon$ and a feasible
$\widehat w$ has empirical objective within $\delta\geq0$ of the simplex
minimum, then
\begin{equation}
 R(\widehat w)-R(w^*)\leq2\epsilon+\delta,
 \label{eq:bound}
\end{equation}
where $w^*$ minimizes population squared relative risk over the simplex.
\end{proposition}
Here $\epsilon$ bounds matrix estimation error and $\delta$ bounds numerical
optimization error. For every feasible $w$, nonnegative unit sum weights give
$|w^\top(\widehat C-C)w|\leq\epsilon\sum_{m,n}w_mw_n=\epsilon$.
Applying this once at $\widehat w$ and once at $w^*$, with empirical near
optimality between them, proves the result. The bound separates estimation
from optimization error but does not certify small $\epsilon$ from five
fields. Standard concentration arguments \citep{hoeffding1963} require
bounded error products and independent fitting fields, assumptions not
established here. Spatial pixels are not independent fitting examples.
The companion package retains the longer concentration calculation, whose
bound is weaker than a trivial bounded risk guarantee at the reported budgets.
Neither calculation establishes near optimality for the unsquared metric $L$.

\clearpage
\subsection{Relationship to the closest prior systems}
\label{app:prior}
Model weighting, automated surrogate selection, and multifidelity field
emulation are established ideas. Table~\ref{tab:prior} identifies these
precedents and the scope of our integration. The contribution is the adapted
portfolio, common parameter input interface, and experiments with counted
fitting examples. The table compares formulations rather than measured
performance. Direct comparisons with complete engineering automation systems
would require matched training, tuning, and information budgets.

\begin{table}[h!]
\centering\small\setlength{\tabcolsep}{3pt}
\caption{\textbf{Prior systems and the scope of this study.} These workflows
were not all evaluated on the present tasks. Adapted baseline implementations
do not substitute for complete system comparisons.}
\label{tab:prior}
\begin{tabular}{p{.24\linewidth}p{.31\linewidth}p{.36\linewidth}}
\toprule
Predecessor & Established contribution & Present investigation \\
\midrule
Engineering ensembles \citep{acar2009,viana2009} & Validation based weighting
and model selection. & Neural field library with several fidelity strategies
and counted fitting examples. \\
Functional output benchmark \citep{brunel2025} & Unified comparison of
multifidelity field emulators. & Heterogeneous neural and reduced basis
candidates combined from saved fields. \\
Stacking and AutoGluon \citep{wolpert1992,erickson2020} & Automated library
combination. & Whole field loss over a fixed portfolio of fidelity strategies. \\
EL MFS and AQBMF \citep{wang2024elmfs,lee2026aqbmf} & Adaptive LF fusion and
MF strategy selection. & Diverse field architectures with parameter only inference. \\
MAESTRO \citep{kim2026maestro} & LF and discrepancy families with screening
and leave one out selection. & Direct, correction, and fine only field models. \\
LANE SI \citep{borisova2024} & Neural and climatological spatial forecast
ensembles. & Parametric emulation with supplied multifidelity simulations. \\
HyperNOs \citep{ghiotto2025} & Automated operator training and hyperparameter
search. & Fitting weights across different fidelity use strategies. \\
MFRNP and MF DeepONet \citep{niu2024,howard2023} & Multifidelity field
architectures. & Common ensemble fitting over multiple adapted families. \\
\bottomrule
\end{tabular}
\end{table}
